\section{Related Work}
\label{sec:related}

\paragraph{LLM unlearning algorithms and benchmarks.}
Modern LLM unlearning uses gradient-based or preference-based fine-tuning to push the model to forget targeted data while preserving the rest. Representative algorithms include \emph{Gradient Ascent} (GradAscent) and gradient-difference variants \citep{jang2023knowledge}, \emph{NPO} (Negative Preference Optimisation) \citep{zhang2024negative}, and \emph{representation misdirection} \citep{li2024wmdp}. Benchmarks such as TOFU \citep{maini2024tofu}, MUSE \citep{shi2024muse}, and WMDP \citep{li2024wmdp} fix specific forget splits and report aggregate forget-quality and utility scores. The unifying \emph{OpenUnlearning} framework \citep{dorna2025openunlearning} integrates $13$ algorithms and $16$ metrics under one Hydra-configured trainer; we vendor it as our stage-2 implementation. These works concentrate on \emph{algorithm choice} given a benchmark forget split. Our work is orthogonal: we hold the algorithm and hyperparameters fixed and study how the choice of $\Df$ itself reshapes side-effects, then turn the observation into a pre-screen for $\Df$.

\paragraph{Side-effects, locality, and knowledge holes.}
Unlearning is increasingly recognised as collateral-damaging. \citet{ko2025probing} introduce \emph{knowledge-hole probing}: standard utility benchmarks miss adjacent and latent knowledge that has been silently corrupted by unlearning, and dedicated probes can recover this damage post-hoc. Studies of model editing \citep{meng2022locating, mitchell2022memory} and influence functions \citep{koh2017understanding, grosse2023studying} show analogous locality / non-locality effects in the gradient-update regime. Our three-layer view operationalises Ko et al.'s phenomenon as a quantitative, dataset-driven measurement: $\Ltwo$ corresponds to their \emph{adjacent} knowledge (same-cluster neighbours), and $\Lthree$ extends the picture to \emph{cross-cluster spillover} that prior probes did not target. Crucially, while their probes act \emph{after} unlearning, our audit acts \emph{before} --- predicting the corruption profile from $\Df$ alone.

\paragraph{Data-first prediction of model behaviour.}
A separate line of work predicts downstream model behaviour from properties of the training data, rather than from running the model. \citet{dang2024curious} predict the adversarial robustness of fine-tuned text transformers from features of the fine-tuning dataset, achieving $30\!\times\!\text{--}\!193\times$ acceleration over the standard model-first attack-then-evaluate pipeline. Related efforts use coreset and data-attribution descriptors \citep{koh2017understanding, ilyas2022datamodels}, dataset diversity statistics \citep{wang2024data}, and curvature of the loss landscape \citep{yao2020pyhessian} to predict generalisation or vulnerability. We carry this paradigm into LLM unlearning: instead of running unlearning and then measuring side-effects, we measure intrinsic geometry of $\Df$ and predict the three-layer corruption profile. To our knowledge, this is the first \emph{data-first} predictor in the unlearning setting.

\paragraph{Forget-set selection and curation.}
A small body of work studies which data is hardest to remove or causes the most damage when removed \citep{thudi2022unrolling, kurmanji2024towards}. The closest setting is influence-based forget-set selection \citep{koh2017understanding}, but influence functions require gradients of the trained model, scale poorly to $8$B-parameter LLMs, and yield a per-example score rather than a $\Df$-level corruption profile. Our auditor uses only $50$ forget-set embeddings and $12$ scalar features; it is two orders of magnitude cheaper and produces a three-layer profile.

\paragraph{Position.}
We sit at the intersection of unlearning side-effect characterisation \citep{ko2025probing} and data-first behaviour prediction \citep{dang2024curious}. The first tells us \emph{what to predict} (locality + spillover); the second tells us \emph{how to predict from data alone}. Our contribution is the quantitative bridge: a three-layer corruption metric tied to a forget-set-side geometric audit.
